% Gauss, Werke X/1 (1917), native85–86, printed78–79.
\SourceLeaf{085}\NoSourceHead
\HistoricTitle{[ZUR THEORIE DER FORMEN.]}
\SourceGroup{I}
\HistoricTitle{[ÜBER POLYGONALZAHLEN.]}
\HistoricCitation{[Eintragung im \textsc{Leiste}*), S. 68.]}
\NumberedItem{1}
\par\indent Die Reihen der Polygonalzahlen lassen sich auch rückwerts fortsetzen,
und dadurch kommt man bei allen, die Trigonalzahlen und Quadratzahlen
ausgenommen, auf Zahlen, welche von den übrigen ganz verschieden sind.
\begin{center}Pentagonalzahlen\end{center}
\[
100,\ 77,\ 57,\ 40,\ 26,\ 15,\ 7,\ 2,\ 0,\ 1,\ 5,\ 12,\ 22,\ 35,\ 51,\ 70,\ 92,\ 117\ldots
\]
\begin{center}Hexagonalzahlen\end{center}
\[
105,\ 78,\ 55,\ 36,\ 21,\ 10,\ 3,\ 0,\ 1,\ 6,\ 15,\ 28,\ 45,\ 66\ldots
\]
\NumberedItem{2}
\par\indent Jede Zahl besteht aus drei Trigonalzahlen.
\par\noindent\hspace*{\parindent}\begin{tabular}{@{}lclcl@{}}
Jede Zahl der Form&$8n+1$&drei&$\square$&$i,pp,pp;\ i,pi,pi$ —\\
&$8n+3$&drei&$\square$&$i,i,i$ —\\
&$8n+5$&drei&$\square$&$i,pp,pi$ —\\
&$8n+7$&vier&$\square$&$i,i,i,pi$ —
\end{tabular}
\par\vfill\noindent\rule{.25\linewidth}{.4pt}\par\smallskip
\begingroup\footnotesize
\noindent[*) Gemeint ist das durchschossene Exemplar von \textsc{Chr. Leiste}, \textit{Die Arithmetik und Algebra}, Wolfenbüttel 1790, das \textsc{Gauss} besessen und dessen Durchschußblätter er von 1796 an bis über 1797 hinaus zu wissenschaftlichen Aufzeichnungen benutzt hat. Die Seitenzahl bezieht sich auf die Seite des Leistetextes, der gegenüber sich die \textsc{Gauss}sche Aufzeichnung befindet.]\par
\endgroup\clearpage
\SourceLeaf{086}\ArithmeticHead{79}{R}{ÜBER POLYGONALZAHLEN.}
\begin{center}\textbf{BEMERKUNG.}\end{center}
\begingroup\small
\par\indent Die vorstehende Aufzeichnung bezieht sich auf den S. 68 stehenden Text des \textsc{Leiste}schen Werkes,
der von den Polygonalzahlen handelt. Solche zum Text gehörige Bemerkungen finden sich auf den Durchschußblättern nur sehr spärlich; wir haben außer der hier abgedruckten noch zu den Seiten 9, 36, 37, 44,
105, 108 welche bemerkt.

Vergl. zu der vorstehenden Aufzeichnung die Tagebuchaufzeichnungen Nr. 17 und 18 vom 3. und
10. Juli 1796 und die darauf bezüglichen Bemerkungen in dem Artikel 12 des \textsc{Bachmann}schen Aufsatzes
»Über \textsc{Gauss}' zahlentheoretische Arbeiten«. Es bedeutet:
\begin{center}\begin{tabular}{rl}
$i$&impar,\\
$pp$&pariter par,\\
$pi$&pariter impar.
\end{tabular}\end{center}
\par\indent Die »series numerorum pentagonalium retro continuata« findet sich schon bei \textsc{Euler} in der Abhandlung \textit{Observatio de summis divisorum}, Novi Comment. Acad. Petrop. 5 (1754/5) 1760, S. 64, Opera omnia,
ser. I, vol. 2, S. 380; er gibt die Zahlen 2, 7, 15, 26, 40, 57, 77, 100, 126.
\begin{flushright}\textsc{Schlesinger.}\end{flushright}
\ClosingRule\clearpage\endgroup
